\chapter{Portfolio Trades and Risk Bids}\label{ch:mx:portfolio-trades-and-risk-bids}

A pension fund asks three brokers for a price at which they will take on a basket of five hundred stocks at tonight's close, knowing only its size,
its sectors and its liquidity. The broker who wins learns the names afterwards, and learns whether it bid too much. This chapter prices such a bid on
a simulated basket: the cost of liquidating it, the risk carried while doing so, and the shading that competition forces on a bidder who does not want
to win only when it is wrong; it then compares a disclosed basket with a blind one and ends with a transition between two portfolios.

\section{Agency and principal portfolio trades}

\begin{definition}[Agency portfolio trade, risk bid]\label{def:mx:portfolio-trades-and-risk-bids:agency}
An \emph{agency portfolio trade}\index{agency portfolio trade} executes a list of orders for a client as agent, the client bearing the costs and the
risk and paying a commission. A \emph{risk bid}\index{risk bid} is a principal price for the whole list: the broker buys (or sells) every position at
an agreed benchmark, usually the close, less a discount, and bears the cost and risk of unwinding them.
\end{definition}

\begin{definition}[Blind risk bid]\label{def:mx:portfolio-trades-and-risk-bids:blind}
A \emph{blind risk bid}\index{blind risk bid} is a \omterm{def:mx:portfolio-trades-and-risk-bids:agency}{risk bid} made on a description of the basket (its size, sectors, liquidity and risk
characteristics) without the names, which are revealed only to the winner once the trade is agreed.
\end{definition}

A portfolio trade (One Quant Book 2, chapter~22) moves a whole list at once: an index rebalance, a fund's inflows, a change of manager. The client
chooses between paying the costs itself (agency) and paying a broker to take them (principal); the \omterm{def:mx:portfolio-trades-and-risk-bids:agency}{risk bid} is priced like any risk transfer
(chapter~21), with two twists: the list is large and diversified, so its risk can be hedged in part, and in the blind version the bidder prices a
description. Kavajecz and Keim (2005) studied such blind principal auctions and found that they gave the liquidity demander a transaction-cost saving
relative to more traditional trading, and liquidity suppliers an efficient way to obtain order flow.

The simulated basket comes from Book~7's synthetic market (\texttt{firm.synthmkt}): 786 names listed on its last day, with a median daily volume of
USD~37.5 million and a median daily volatility of 171 basis points; half-spreads follow a size rule (an assumption: a median full spread of 5.9 basis
points), and a ten-factor statistical risk model on 250 days of returns gives the covariance (Book~7's \texttt{firm.riskmodel}). The pension fund's
basket holds 500 of them, USD~250 million, with positions from USD~38\,000 to USD~4.8 million, a median of 1.0\% of a day's volume.

\section{The cost of liquidating a basket}

\begin{definition}[Basket liquidity profile]\label{def:mx:portfolio-trades-and-risk-bids:profile}
A \emph{basket liquidity profile}\index{basket liquidity profile} summarises a basket without its names: the share of its notional in buckets of
position size relative to daily volume, by sector, with its beta, volatility and other risk characteristics.
\end{definition}

Each position is sold at no more than 20\% of its daily volume, so a position of $k$ days' volume takes $\max(1,5k)$ days; its cost is its
half-spread plus the square-root law's $Y\sigma\sqrt{\min(Q/V,0.2)}$ (chapter~11, $Y=0.7$). The basket's liquidity profile is
\cref{fig:mx:portfolio-trades-and-risk-bids:profile}: 98.8\% of its notional is below 5\% of a day's volume, and no position needs more than a day. The
winner's cost of unwinding it is 14.9 basis points of the basket (\texttt{firm\_riskbid.liquidation}).

\begin{omfigure}
\begin{tikzpicture}
\begin{axis}[width=9.4cm,height=5.2cm,ybar,bar width=14pt,ylabel={\small share of the basket's notional (\%)},
  xtick={0,1,2,3,4,5},xticklabels={{$<0.5\%$},{0.5--1\%},{1--2\%},{2--5\%},{5--10\%},{10--20\%}},
  xlabel={\small position as a share of daily volume},tick label style={font=\footnotesize},ymin=0,ymax=40,
  grid=major,grid style={gray!25},enlarge x limits=0.12,table/col sep=comma]
\addplot[fill=omDef!70,draw=omDef] table[x=i,y=share]{figdata/microstructure/22-portfolio-trades-and-risk-bids/profile.csv};
\end{axis}
\end{tikzpicture}

\omcaption{The basket's liquidity profile: its notional by position size relative to daily volume, what a blind bidder is shown (with the sector
weights). Data: \texttt{mx\_riskbid.blind}.}\label{fig:mx:portfolio-trades-and-risk-bids:profile}
\end{omfigure}

\section{Pricing a risk bid}

While it sells, the winner holds what is left. With each position sold down in a straight line over its days $d_i$, the variance of the P\&L is
$\sum_{ij}w_iw_jC_{ij}I(d_i,d_j)$, with $C$ the daily covariance and $I(a,b)=a/2-a^2/(6b)$ for $a\le b$ the overlap of the two holdings (for one day,
a third of a day's variance): this is the multi-asset Almgren--Chriss risk term (Almgren and Chriss, 2001; chapter~14) for straight-line schedules. For
the basket, one standard deviation is 40.4 basis points; sold short in index futures at once to remove its market beta, 12.5.

\omcode{code/firm/riskbid/firm_riskbid.py}{43}{57}{The risk of holding a basket while it is sold down: each pair of positions contributes its covariance times the time their straight-line holdings overlap; a market hedge removes the factor's part of the covariance.}\label{lst:mx:portfolio-trades-and-risk-bids:risk}

The bid, as a discount to the close, adds four parts: the expected cost, a charge for the risk (here one standard deviation of the hedged P\&L, the
desk's choice), the shading that competition requires (the next section), and a margin (2 basis points). For the disclosed basket: 14.9 of cost,
12.5 of risk, 4.2 of shading and 2.0 of margin, a bid of 33.6 basis points below the close (\cref{fig:mx:portfolio-trades-and-risk-bids:bids}).

\section{Blind bids and the winner's curse}

Every bidder estimates the same cost with its own error; the one whose estimate is lowest wins, and it wins most often when its estimate is too low.
This is the winner's curse of a common-value auction (One Quant Book 2, chapter~30). If each of four bidders (the desk and three competitors) estimates
the cost with an independent error of standard deviation $s$, the winner's estimate is on average $s\,\E[\max\text{ of 4 standard normals}]=1.03s$
too low, and a bidder who does not shade by that amount loses it on the trades it wins (\texttt{firm\_riskbid.curse}).

\omcode{code/firm/riskbid/firm_riskbid.py}{69}{80}{The winner's curse and the auction that checks it: the shading is the expected lead of the best of n + 1 estimates; the simulated auction returns the winner's average profit.}\label{lst:mx:portfolio-trades-and-risk-bids:curse}

With estimation errors of 15\% of the cost and risk (4.1 basis points for the disclosed basket), the shading is 4.2 basis points; in 200\,000
simulated auctions, bidders who add only the 2-basis-point margin make $-2.2$ basis points on the trades they win, and bidders who also shade make the
2.0 they wanted.

A blind bidder sees only the profile. To price it, the desk draws 300 baskets with the same profile (each position keeps its sector, its notional and
its bucket of days of volume, with a name drawn from those that fit) and prices each: their cost averages 17.0 basis points (standard deviation 0.4)
and their hedged risk 13.0 (1.1). The profile does not reveal that the fund held more of the liquid names than a random draw would: the blind estimate
is 2.1 basis points above the disclosed cost. The larger estimation error adds 0.2 to the shading. The blind bid is 36.4 basis points, 2.8 more than the
disclosed one; the client pays the difference for not showing its names before the trade.

\begin{omfigure}
\begin{tikzpicture}
\begin{axis}[width=10.4cm,height=4.2cm,xbar stacked,bar width=13pt,xlabel={\small bid below the close (basis points)},
  ytick={0,1},yticklabels={blind,disclosed},tick label style={font=\footnotesize},xmin=0,xmax=40,ymin=-0.6,ymax=1.6,
  legend columns=4,legend style={font=\footnotesize,at={(0.5,-0.42)},anchor=north,draw=none,fill=none,
  /tikz/every even column/.append style={column sep=6pt}},table/col sep=comma]
\addplot[fill=omDef!70,draw=omDef] table[y=y,x=cost]{figdata/microstructure/22-portfolio-trades-and-risk-bids/bids.csv};
\addplot[fill=omThm!60,draw=omThm] table[y=y,x=risk]{figdata/microstructure/22-portfolio-trades-and-risk-bids/bids.csv};
\addplot[fill=omProp!70,draw=omProp] table[y=y,x=curse]{figdata/microstructure/22-portfolio-trades-and-risk-bids/bids.csv};
\addplot[fill=gray!40,draw=gray] table[y=y,x=margin]{figdata/microstructure/22-portfolio-trades-and-risk-bids/bids.csv};
\legend{cost,risk charge,winner's curse,margin}
\end{axis}
\end{tikzpicture}

\omcaption{The disclosed and blind bids on the same basket, part by part (basis points of the basket below the close): the blind bidder's cost and risk
are its estimates from the liquidity profile. Data: \texttt{mx\_riskbid.bids}.}\label{fig:mx:portfolio-trades-and-risk-bids:bids}
\end{omfigure}

The blind bid still has a use: the client's names do not leak to the losers, who might otherwise trade ahead of the winner's unwind (chapter~9's
leakage). Whether the 2.8 basis points buy enough protection depends on how many dealers would see the names and what they would do with them.

\section{Transitions}

A transition moves a fund from one portfolio (a manager it is leaving) to another. The transition manager (One Quant Book 8, chapter~29) nets the names
the two portfolios share, crosses what it can internally against other flow, and trades the rest. For two portfolios of USD~250 million with 150 names
in common: selling one and buying the other outright trades USD~500 million and costs 50.4 basis points of the fund; netting the shared names trades
USD~322 million at 32.8; crossing 20\% of the remainder internally leaves USD~258 million to trade, at 24.0 (\texttt{firm\_riskbid.transition}).
Internal crossing (One Quant Book 7, chapter~27) and the central risk book (Book~9, chapter~25) are the same idea at the scale of a bank: the cheapest
trade is the one that never reaches the market.

\section{Tutorial: pricing a risk bid}\label{tut:mx:portfolio-trades-and-risk-bids:tut}

\begin{tutorial}
\textbf{Goal.} Price a principal bid for a 500-name basket, disclosed and blind, with its winner's-curse shading, and cost a transition.
\textbf{End state:} \cref{fig:mx:portfolio-trades-and-risk-bids:profile,fig:mx:portfolio-trades-and-risk-bids:bids} and the numbers of sections 2 to 5.
\begin{tutsteps}
\item \textbf{Universe.} \texttt{mx\_riskbid.universe()} on \texttt{firm.synthmkt} and \texttt{firm.riskmodel.pca\_model}; \texttt{basket()}.
\item \textbf{Cost and risk.} \texttt{firm\_riskbid.liquidation}, \texttt{risk\_sd}; \texttt{price(names, notional)}.
\item \textbf{Blind.} \texttt{profile}; \texttt{blind()}.
\item \textbf{Bid.} \texttt{curse}, \texttt{auction}; \texttt{bids()}; \texttt{transition\_study()}; draw with \texttt{fig\_riskbid.py}.
\end{tutsteps}
\textbf{What to change next.} Add \omterm{def:mx:metaorders-latent-liquidity-and-cross-impact:cross}{cross-impact} between the names (chapter~13); let the competitors' errors be correlated (a shared model); make
the basket less liquid than its profile suggests.
\end{tutorial}

\section{Build: risk bids}\label{bld:mx:portfolio-trades-and-risk-bids:riskbid}

\begin{build}
\bfield{Purpose} The pricing of principal baskets and transitions for the desk of chapter~20 and the bank desks of Book~9.
\bfield{Interface} \texttt{liquidation(notional, adv, sigma, spread\_bp, cap, y)}, \texttt{risk\_sd(notional, cov, days, hedge)},
\texttt{profile(notional, adv, sector, edges)}, \texttt{curse(sigma\_est, competitors)}, \texttt{auction(true\_cost, sigma\_est, bidders, shade,
margin, trials, seed)}, \texttt{transition(old, new, cross\_share)}.
\bfield{Rules} Costs in basis points of the basket's gross notional; straight-line liquidation within each name's days; a hedge removes a factor at
once; bidders' errors independent.
\bfield{Acceptance tests} \texttt{code/firm/riskbid/tests/}: the liquidation rule; the profile; the straight-line risk against a numerical integral
and a perfect hedge; the curse's shading makes the winner's profit the margin; the transition's netting and crossing.
\bfield{Stretch} \omterm{def:mx:metaorders-latent-liquidity-and-cross-impact:cross}{Cross-impact}; correlated bidders; optimal liquidation schedules instead of straight lines.
\end{build}

\begin{omsources}
\item R.~Almgren and N.~Chriss, ``Optimal execution of portfolio transactions'', \emph{Journal of Risk} 3(2), 2001.
\item K.~A.~Kavajecz and D.~B.~Keim, ``Packaging liquidity: blind auctions and transaction efficiencies'', \emph{Journal of Financial and
Quantitative Analysis} 40(3), 2005.
\end{omsources}

\section{Exercises}

\begin{exercise}[$\star$]\label{exo:mx:portfolio-trades-and-risk-bids:1}
A position of 30\% of a day's volume in a stock with a daily volatility of 200 basis points and a 6-basis-point spread: how many days does it take at
a 20\% cap, and what does it cost?
\end{exercise}

\begin{exercise}[$\star$]\label{exo:mx:portfolio-trades-and-risk-bids:2}
One position sold in a straight line over three days with a daily volatility of 150 basis points: what is the standard deviation of its P\&L, in
basis points of the position?
\end{exercise}

\begin{exercise}[$\star$]\label{exo:mx:portfolio-trades-and-risk-bids:3}
What shading does a bidder need against five competitors with an estimation error of 4 basis points?
\end{exercise}

\begin{exercise}[$\star\star$]\label{exo:mx:portfolio-trades-and-risk-bids:4}
Why does hedging with index futures cut the basket's risk from 40.4 to 12.5 basis points, and why not further?
\end{exercise}

\begin{exercise}[$\star\star$]\label{exo:mx:portfolio-trades-and-risk-bids:5}
Why does the blind bidder overestimate this basket's cost, and could it underestimate another's?
\end{exercise}

\begin{exercise}[$\star\star$]\label{exo:mx:portfolio-trades-and-risk-bids:6}
If all four bidders use the same vendor's cost model, what happens to the winner's curse?
\end{exercise}

\begin{exercise}[$\star\star\star$]\label{exo:mx:portfolio-trades-and-risk-bids:7}
\emph{Coding.} Price the disclosed basket with a risk charge of two standard deviations instead of one. What is the bid?
\end{exercise}

\begin{exercise}[$\star\star\star$]\label{exo:mx:portfolio-trades-and-risk-bids:8}
\emph{Find the flaw.} ``We won 40\% of the \omterm{def:mx:portfolio-trades-and-risk-bids:agency}{risk bids} we made last year, so our pricing is about right.''
\end{exercise}

\section{Problem: Pricing the Basket Blind}

\begin{problem}[{Weekend problem --- pricing the basket blind}]\label{pb:mx:portfolio-trades-and-risk-bids:1}
A pension fund offers a 500-name basket for a principal bid at the close; your desk competes with three others.

\medskip\noindent\textbf{Part I --- The trade.}
\begin{enumerate}
\item Distinguish agency and principal portfolio trades.
\item Define a \omterm{def:mx:portfolio-trades-and-risk-bids:agency}{risk bid} and a \omterm{def:mx:portfolio-trades-and-risk-bids:blind}{blind risk bid}.
\item Describe the simulated universe and the basket.
\item What did Kavajecz and Keim find about blind principal auctions?
\end{enumerate}

\medskip\noindent\textbf{Part II --- Cost and risk.}
\begin{enumerate}[resume]
\item State the liquidation rule and give the basket's cost.
\item Derive the overlap $I(a,b)$ of two straight-line holdings.
\item Give the basket's risk with and without the market hedge.
\item What does the liquidity profile show?
\end{enumerate}

\medskip\noindent\textbf{Part III --- Competition.}
\begin{enumerate}[resume]
\item Explain the winner's curse and derive the shading.
\item What do bidders make without shading and with it?
\item How does the blind bidder estimate cost and risk, and what does it find?
\item Give the disclosed and blind bids part by part.
\item \emph{State the named result}: the bid in basis points for the disclosed and the blind basket, and the winner's-curse adjustment against three
competitors.
\end{enumerate}

\medskip\noindent\textbf{Part IV --- Transitions and judgement.}
\begin{enumerate}[resume]
\item Cost the transition outright, netted and with internal crossing.
\item Why is internal crossing cheaper than any execution?
\item When is a blind bid worth its premium to the client?
\item Which assumption of the pricing would you check first?
\item How would you price the bid if competitors shared your model?
\item What would you tell the fund about agency against principal?
\item In one sentence: what does a \omterm{def:mx:portfolio-trades-and-risk-bids:agency}{risk bid} sell?
\end{enumerate}
\end{problem}

\section{Interview questions}

\begin{interviewq}[$\star$ \iqroles{trader}]\label{iq:mx:portfolio-trades-and-risk-bids:1}
A client asks for a risk price on a basket at the close. What do you need to know, and how do you build the price?
\end{interviewq}

\begin{interviewq}[$\star\star$ \iqroles{researcher}]\label{iq:mx:portfolio-trades-and-risk-bids:2}
Explain the winner's curse in a risk-bid auction and how you would correct for it.
\end{interviewq}

\begin{interviewq}[$\star\star$ \iqroles{risk}]\label{iq:mx:portfolio-trades-and-risk-bids:3}
How would you measure the risk of holding a basket while you unwind it, and what can you hedge?
\end{interviewq}

\begin{interviewq}[$\star\star$ \iqroles{bank}]\label{iq:mx:portfolio-trades-and-risk-bids:4}
Your desk keeps losing money on the blind bids it wins. What do you check?
\end{interviewq}

\begin{interviewq}[$\star\star$ \iqroles{trader}]\label{iq:mx:portfolio-trades-and-risk-bids:5}
Plan a transition from one equity manager to another for a pension fund.
\end{interviewq}

\begin{interviewq}[$\star\star\star$ \iqroles{researcher}]\label{iq:mx:portfolio-trades-and-risk-bids:6}
How would you estimate the cost of a basket you cannot see from its liquidity profile, and how would you know if the estimate was biased?
\end{interviewq}
